\documentclass[12pt, a4paper]{article}

\usepackage[utf8]{inputenc}
\usepackage[T1]{fontenc}
\usepackage[brazil]{babel}
\usepackage{amsmath, amssymb}
\usepackage{geometry}
\usepackage{booktabs}
\usepackage{multirow}
\usepackage{graphicx}
\usepackage{xcolor}
\usepackage{hyperref}
\usepackage{fancyhdr}
\usepackage{titlesec}
\usepackage{enumitem}
\usepackage{float}
\usepackage{caption}
\usepackage{longtable}
\usepackage{colortbl}

\geometry{left=2.5cm, right=2.5cm, top=3cm, bottom=3cm}

% Cores
\definecolor{azulescuro}{RGB}{0, 51, 102}
\definecolor{verdepositivo}{RGB}{0, 128, 0}
\definecolor{vermelhoneg}{RGB}{178, 34, 34}
\definecolor{amareloalerta}{RGB}{204, 153, 0}
\definecolor{cinzaclaro}{RGB}{240, 240, 240}
\definecolor{cinzatabela}{RGB}{245, 245, 245}

% Configurações
\hypersetup{
    colorlinks=true,
    linkcolor=azulescuro,
    urlcolor=azulescuro,
    citecolor=azulescuro
}

\pagestyle{fancy}
\fancyhf{}
\fancyhead[L]{\small\color{azulescuro}Análise Quantitativa --- Track\&Field (TFCO4.SA)}
\fancyhead[R]{\small\color{azulescuro}Março 2026}
\fancyfoot[C]{\thepage}
\renewcommand{\headrulewidth}{0.4pt}

\titleformat{\section}{\large\bfseries\color{azulescuro}}{}{0em}{}[\titlerule]
\titleformat{\subsection}{\normalsize\bfseries\color{azulescuro}}{}{0em}{}

\begin{document}

% ============================================================================
% CAPA
% ============================================================================
\begin{titlepage}
\centering
\vspace*{3cm}

{\Huge\bfseries\color{azulescuro} Análise Quantitativa}\\[0.5cm]
{\LARGE\color{azulescuro} Track\&Field (TFCO4.SA)}\\[1cm]

\rule{\textwidth}{1pt}\\[0.5cm]

{\large Valuation, Qualidade, Risco e Otimização}\\[0.3cm]
{\large Comparação com Peers do Setor de Varejo/Moda}\\[2cm]

\begin{tabular}{ll}
\textbf{Data da Análise:} & 03 de Março de 2026 \\
\textbf{Período de Dados:} & Mar/2021 -- Mar/2026 (5 anos) \\
\textbf{Ativo Principal:} & TFCO4.SA (B3) \\
\textbf{Setor:} & Varejo / Moda Esportiva \& Lifestyle \\
\textbf{Taxa Livre de Risco:} & Selic 14,25\% a.a. \\
\end{tabular}

\vfill

{\small\textit{Nota: Esta análise NÃO constitui recomendação de compra ou venda de ativos.\\
É um framework quantitativo para apoio à decisão de investimento.}}

\end{titlepage}

\tableofcontents
\newpage

% ============================================================================
% 1. EXECUTIVE SUMMARY
% ============================================================================
\section{Executive Summary}

\subsection{Resultado da Otimização}

O framework quantitativo avaliou 6 ativos do setor de varejo/moda na B3 utilizando critérios de \textbf{Valuation}, \textbf{Qualidade}, \textbf{Risco} e \textbf{Monte Carlo}, com otimização via \textit{Simulated Annealing} (inspirado em computação quântica).

\begin{itemize}[leftmargin=*]
    \item \textbf{Ativo selecionado (SA/QUBO):} Grendene (GRND3.SA) --- Score: 0,7289
    \item \textbf{Ranking completo:} Grendene $>$ Vivara $>$ Track\&Field $>$ Grupo SBF $>$ Lojas Renner $>$ Alpargatas
    \item \textbf{Carteira ótima (scipy):} Grendene 40\% + Vivara 40\% + Track\&Field 20\%
\end{itemize}

\subsection{Destaques Track\&Field}

\begin{center}
\begin{tabular}{lrl}
\toprule
\textbf{Métrica} & \textbf{Valor} & \textbf{Comentário} \\
\midrule
Preço atual & R\$ 16,58 & \\
Posição 52 semanas & 80,1\% & Próximo do topo \\
Market Cap & R\$ 2,51 B & Small/mid-cap \\
P/E Trailing & 21,26x & Premium vs peers \\
P/E Forward & 13,37x & Desconto forward \\
EV/EBITDA & 13,26x & Acima da mediana do setor \\
ROE & 28,27\% & \textcolor{verdepositivo}{Superior aos peers (16,1\%)} \\
Retorno 1 ano & +72,2\% & \textcolor{verdepositivo}{Alpha +57,9\% vs CDI} \\
Volatilidade & 36,9\% & \textcolor{verdepositivo}{Menor que peers (41,9\%)} \\
Beta (Ibovespa) & 1,003 & Próximo do mercado \\
Max Drawdown & $-48,9\%$ & Risco relevante \\
\bottomrule
\end{tabular}
\end{center}

% ============================================================================
% 2. TESE DE INVESTIMENTO
% ============================================================================
\section{Tese de Investimento}

\subsection{Contexto: Track\&Field}

A Track\&Field é uma das principais marcas brasileiras de moda fitness e lifestyle premium. O modelo de negócio é \textbf{asset-light}, baseado em franquias, e-commerce e lojas próprias, o que reduz a necessidade de capital e melhora o retorno sobre o patrimônio.

\textbf{Catalisadores positivos:}
\begin{enumerate}[leftmargin=*]
    \item \textbf{Tendência secular de wellness:} O mercado de fitness/lifestyle segue em expansão pós-pandemia, favorecendo marcas especializadas
    \item \textbf{Modelo de franquias:} Permite crescimento acelerado com baixo capex próprio (30--50 lojas/ano)
    \item \textbf{Expansão de e-commerce:} Canal D2C (direct-to-consumer) com margens superiores
    \item \textbf{Margens premium:} Margem bruta de 56\% e ROE de 28\% são excepcionais para varejo brasileiro
    \item \textbf{Geração de caixa:} FCF positivo e distribuição consistente de dividendos
\end{enumerate}

\textbf{Riscos principais:}
\begin{enumerate}[leftmargin=*]
    \item Alta sensibilidade ao ciclo econômico (Selic, emprego, confiança do consumidor)
    \item Setor de varejo de moda é competitivo e cíclico
    \item Pressão de custo por importação de matéria-prima (exposição cambial)
    \item Small/mid-cap com liquidez mais restrita na B3
    \item Dependência de shopping centers
    \item Risco de execução na expansão acelerada
\end{enumerate}

% ============================================================================
% 3. VALUATION COMPARATIVO
% ============================================================================
\section{Valuation Comparativo}

\subsection{Múltiplos de Mercado}

\begin{center}
\small
\begin{tabular}{lccccccc}
\toprule
\textbf{Empresa} & \textbf{E. Yield} & \textbf{FCF Yield} & \textbf{EV/EBITDA} & \textbf{P/E} & \textbf{P/B} & \textbf{Div. Yield} & \textbf{Cresc. Rec.} \\
\midrule
\rowcolor{cinzatabela}
Track\&Field & 0,05 & 0,01 & 13,26x & 21,26x & 4,74x & 1,6\%$^*$ & 31\% \\
Alpargatas & 0,04 & 0,02 & 23,29x & 26,86x & 2,16x & 4,3\% & 7\% \\
\rowcolor{cinzatabela}
Grupo SBF & 0,12 & 0,22 & 7,58x & 8,49x & 0,86x & 4,3\% & 9\% \\
Lojas Renner & 0,09 & 0,11 & 7,25x & 11,16x & 1,46x & 5,9\% & 5\% \\
\rowcolor{cinzatabela}
Grendene & 0,16 & 0,04 & 9,49x & 6,16x & 1,10x & 27,1\%$^*$ & 1\% \\
Vivara & 0,10 & $-0,05$ & 10,44x & 9,79x & 2,47x & 4,4\% & 18\% \\
\bottomrule
\end{tabular}
\end{center}
{\footnotesize $^*$ Dividend Yield possivelmente distorcido por dados da API yfinance.}

\subsection{Análise}

A Track\&Field negocia a \textbf{P/E de 21,26x} (trailing), acima da mediana do grupo (10,28x), refletindo o prêmio por crescimento. No entanto, o \textbf{P/E forward de 13,37x} sugere expectativa de forte expansão de lucros, mais alinhado com o setor.

O \textbf{EV/EBITDA de 13,26x} é o segundo mais alto do grupo (atrás apenas de Alpargatas a 23,29x), mas é justificado pelo crescimento de receita de 31\% --- significativamente acima dos peers.

\textbf{Valuation Score (Z-Score):} $-0,93$, indicando que em termos puramente de múltiplos, a Track\&Field está \textit{relativamente cara} frente aos peers. O mercado já precifica parte do crescimento.

% ============================================================================
% 4. QUALIDADE
% ============================================================================
\section{Métricas de Qualidade}

\subsection{Margens e Retornos}

\begin{center}
\small
\begin{tabular}{lcccccc}
\toprule
\textbf{Empresa} & \textbf{M. Bruta} & \textbf{M. Oper.} & \textbf{M. Líq.} & \textbf{ROE} & \textbf{ROA} & \textbf{FCF/Rev} \\
\midrule
\rowcolor{cinzatabela}
Track\&Field & \textbf{56\%} & \textbf{18\%} & 14\% & \textcolor{verdepositivo}{\textbf{28,3\%}} & \textbf{16\%} & 4\% \\
Alpargatas & 48\% & 16\% & 8\% & 9\% & 3\% & 5\% \\
\rowcolor{cinzatabela}
Grupo SBF & 48\% & 7\% & 4\% & 11\% & 4\% & 8\% \\
Lojas Renner & 61\% & 8\% & 9\% & 14\% & 6\% & 10\% \\
\rowcolor{cinzatabela}
Grendene & 47\% & 9\% & \textbf{27\%} & 18\% & 5\% & 7\% \\
Vivara & \textbf{72\%} & \textbf{32\%} & 26\% & \textcolor{verdepositivo}{\textbf{28,0\%}} & 12\% & $-14\%$ \\
\bottomrule
\end{tabular}
\end{center}

\subsection{Alavancagem e Liquidez}

\begin{center}
\begin{tabular}{lccc}
\toprule
\textbf{Empresa} & \textbf{D/E} & \textbf{Dív. Líq./EBITDA} & \textbf{Líq. Corrente} \\
\midrule
\rowcolor{cinzatabela}
Track\&Field & 0,32x & 0,71x & 2,86x \\
Alpargatas & 0,19x & $-0,50x$ & 2,40x \\
\rowcolor{cinzatabela}
Grupo SBF & 1,00x & 3,43x & 1,44x \\
Lojas Renner & 0,28x & 0,54x & 1,76x \\
\rowcolor{cinzatabela}
Grendene & 0,02x & $-2,69x$ & 6,07x \\
Vivara & 0,43x & 1,31x & 5,17x \\
\bottomrule
\end{tabular}
\end{center}

\subsection{Análise}

\textbf{Hipótese H1 --- CONFIRMADA:} A Track\&Field apresenta o \textbf{maior ROE do grupo} (28,27\% vs média de 16,12\% dos peers), combinado com margens operacionais robustas (18\%) e alavancagem controlada (D/E = 0,32x, Dív.Líq./EBITDA = 0,71x).

O \textbf{Quality Score} de 0,54 (Z-Score) posiciona a Track\&Field como a 3ª melhor em qualidade, atrás de Vivara (1,28) e Grendene (0,51). Destaque para Vivara com margens brutas excepcionais (72\%) e Grendene com posição de caixa líquido ($-2,69x$ dív.líq./EBITDA).

% ============================================================================
% 5. RISCO
% ============================================================================
\section{Métricas de Risco}

\subsection{Retorno, Volatilidade e VaR}

\begin{center}
\small
\begin{tabular}{lcccccc}
\toprule
\textbf{Empresa} & \textbf{Ret. Anual} & \textbf{Vol. Anual} & \textbf{Sharpe} & \textbf{Max DD} & \textbf{VaR 95\%} & \textbf{CVaR 95\%} \\
\midrule
\rowcolor{cinzatabela}
Track\&Field & 6\% & 37\% & $-0,22$ & $-49\%$ & $-3,7\%$ & $-5,1\%$ \\
Alpargatas & $-15\%$ & 45\% & $-0,65$ & $-90\%$ & $-4,4\%$ & $-6,6\%$ \\
\rowcolor{cinzatabela}
Grupo SBF & $-10\%$ & 53\% & $-0,47$ & $-84\%$ & $-5,3\%$ & $-7,2\%$ \\
Lojas Renner & $-11\%$ & 42\% & $-0,61$ & $-71\%$ & $-4,3\%$ & $-6,0\%$ \\
\rowcolor{cinzatabela}
Grendene & 6\% & 31\% & $-0,25$ & $-50\%$ & $-3,0\%$ & $-4,3\%$ \\
Vivara & 9\% & 38\% & $-0,15$ & $-52\%$ & $-3,7\%$ & $-5,3\%$ \\
\bottomrule
\end{tabular}
\end{center}

\textbf{Nota:} Sharpe negativo para todos os ativos reflete o efeito da alta taxa de juros brasileira (Selic 14,25\%). Neste cenário, o prêmio de risco de equities é estruturalmente baixo.

\subsection{Betas (Sensibilidade ao Ibovespa)}

\begin{center}
\begin{tabular}{lcc}
\toprule
\textbf{Empresa} & \textbf{Beta} & \textbf{Interpretação} \\
\midrule
\rowcolor{cinzatabela}
Track\&Field & 1,003 & Neutro (segue o mercado) \\
Alpargatas & 1,293 & Agressivo \\
\rowcolor{cinzatabela}
Grupo SBF/Centauro & 1,735 & Muito agressivo \\
Lojas Renner & 1,543 & Agressivo \\
\rowcolor{cinzatabela}
Grendene & 0,874 & \textcolor{verdepositivo}{Defensivo} \\
Vivara & 1,279 & Agressivo \\
\bottomrule
\end{tabular}
\end{center}

\textbf{Hipótese H3 --- CONFIRMADA:} A Track\&Field apresenta volatilidade anualizada de 36,9\%, inferior à média dos peers (41,9\%). Grendene é a mais defensiva (31\% vol, beta 0,87).

% ============================================================================
% 6. RETORNOS e ANÁLISE DETALHADA
% ============================================================================
\section{Análise Detalhada --- Track\&Field}

\subsection{Retornos Acumulados}

\begin{center}
\begin{tabular}{lcc}
\toprule
\textbf{Período} & \textbf{TFCO4} & \textbf{CDI} \\
\midrule
1 mês & +2,7\% & $\approx$1,1\% \\
3 meses & $-6,5\%$ & $\approx$3,4\% \\
6 meses & +3,9\% & $\approx$6,9\% \\
\textbf{1 ano} & \textcolor{verdepositivo}{\textbf{+72,2\%}} & \textbf{14,3\%} \\
2 anos & +37,2\% & $\approx$27\% \\
\midrule
\textbf{Alpha vs CDI (1 ano)} & \multicolumn{2}{c}{\textcolor{verdepositivo}{\textbf{+57,9 p.p.}}} \\
\bottomrule
\end{tabular}
\end{center}

O retorno de 72,2\% no último ano demonstra forte recuperação e reavaliação pelo mercado. No entanto, o retorno de 2 anos (+37,2\%) mostra que houve períodos de queda significativa antes.

\subsection{Drawdown}

\begin{itemize}[leftmargin=*]
    \item \textbf{Max Drawdown (5 anos):} $-48,9\%$ --- queda severa, típica de small/mid-caps em momentos de estresse
    \item \textbf{Drawdown Atual:} $-9,4\%$ --- ainda abaixo do topo recente, com espaço para recuperação
\end{itemize}

\subsection{Análise Técnica}

\begin{center}
\begin{tabular}{lrl}
\toprule
\textbf{Indicador} & \textbf{Valor} & \textbf{Sinal} \\
\midrule
Preço atual & R\$ 16,58 & \\
SMA 20 & R\$ 16,90 & Abaixo (curto prazo fraco) \\
SMA 50 & R\$ 16,23 & Acima \\
SMA 200 & R\$ 15,93 & Acima (+4,1\%) \\
\midrule
Tendência & \multicolumn{2}{c}{Alta (Golden Cross)} \\
RSI (14 dias) & 41,7 & Neutro \\
\bottomrule
\end{tabular}
\end{center}

A tendência macro é de \textbf{alta} (SMA 50 $>$ SMA 200), apesar de fraqueza de curto prazo (preço abaixo da SMA 20). RSI neutro em 41,7 sugere ausência de sobrecompra.

\subsection{Correlações}

\begin{center}
\begin{tabular}{lc}
\toprule
\textbf{Par} & \textbf{Correlação} \\
\midrule
TFCO4 vs Ibovespa & 0,468 \\
TFCO4 vs Dólar (USDBRL) & $-0,030$ \\
TFCO4 vs Lojas Renner & 0,466 \\
TFCO4 vs Vivara & 0,463 \\
TFCO4 vs Grupo SBF & 0,404 \\
TFCO4 vs Grendene & 0,337 \\
TFCO4 vs Alpargatas & 0,327 \\
\bottomrule
\end{tabular}
\end{center}

Correlação moderada com o Ibovespa (0,47) e quase zero com o dólar, indicando que TFCO4 é primariamente um ativo de \textit{consumo doméstico} com baixa exposição cambial direta.

% ============================================================================
% 7. REGRESSÃO MULTIFATORIAL
% ============================================================================
\section{Regressão Multifatorial}

\subsection{Modelo: $r_{TFCO4} = \alpha + \beta_1 \cdot r_{IBOV} + \beta_2 \cdot r_{USDBRL} + \varepsilon$}

\begin{center}
\begin{tabular}{lccc}
\toprule
\textbf{Empresa} & \textbf{Alpha (anual)} & $\boldsymbol{\beta_{IBOV}}$ & $\boldsymbol{\beta_{USDBRL}}$ \\
\midrule
\rowcolor{cinzatabela}
Track\&Field & $-0,034$ & $1,004^{***}$ & $-0,093$ \\
Alpargatas & $-0,275$ & $1,293^{***}$ & $-0,028$ \\
\rowcolor{cinzatabela}
Grupo SBF & $-0,272$ & $1,737^{***}$ & $-0,174^{**}$ \\
Lojas Renner & $-0,259$ & $1,544^{***}$ & $-0,024$ \\
\rowcolor{cinzatabela}
Grendene & $-0,019$ & $0,875^{***}$ & $-0,041$ \\
Vivara & $-0,037$ & $1,281^{***}$ & $-0,141^{**}$ \\
\bottomrule
\end{tabular}
\end{center}

{\footnotesize $^{***}p < 0,01$; $^{**}p < 0,05$; $^{*}p < 0,10$}

\textbf{Destaques:}
\begin{itemize}[leftmargin=*]
    \item Track\&Field tem \textbf{alpha anualizado quase zero} ($-3,4\%$), indicando retorno próximo do esperado dado o beta
    \item $R^2 = 0,22$: apenas 22\% da variação é explicada pelo mercado e câmbio (idiossincrasia alta)
    \item Grupo SBF tem o maior beta (1,74) e Grendene o menor (0,87)
    \item Sensibilidade cambial ($\beta_{USDBRL}$) é estatisticamente significativa apenas para Grupo SBF e Vivara
\end{itemize}

% ============================================================================
% 8. MONTE CARLO
% ============================================================================
\section{Simulação Monte Carlo (12 meses)}

Simulação com 10.000 trajetórias usando distribuição \textbf{t-Student} ($\nu = 5$) para capturar caudas gordas.

\begin{center}
\small
\begin{tabular}{lccccc}
\toprule
\textbf{Empresa} & \textbf{E[Ret]} & \textbf{Mediana} & \textbf{VaR 95\%} & \textbf{P($>$0)} & \textbf{P($>$CDI)} \\
\midrule
\rowcolor{cinzatabela}
Track\&Field & 18,9\% & 7\% & $-51,4\%$ & 56\% & 43,9\% \\
Alpargatas & 2,8\% & $-14\%$ & $-67,2\%$ & 40\% & 31,6\% \\
\rowcolor{cinzatabela}
Grupo SBF & 14,8\% & $-8\%$ & $-69,1\%$ & 45\% & 36,3\% \\
Lojas Renner & 2,1\% & $-12\%$ & $-64,0\%$ & 41\% & 31,8\% \\
\rowcolor{cinzatabela}
Grendene & 15,6\% & 7\% & $-44,9\%$ & 57\% & 43,5\% \\
Vivara & 23,0\% & 9\% & $-51,0\%$ & 57\% & 46,3\% \\
\bottomrule
\end{tabular}
\end{center}

\textbf{Interpretação para TFCO4:}
\begin{itemize}[leftmargin=*]
    \item Retorno esperado de 18,9\% é atrativo, mas a mediana (7\%) é significativamente menor --- indica assimetria positiva (upside potencial grande, mas cenário provável mais modesto)
    \item VaR 95\% de $-51,4\%$: há 5\% de chance de perder mais de metade do capital em 12 meses
    \item Probabilidade de bater o CDI: apenas 43,9\% --- refletindo o custo de oportunidade da Selic alta
    \item Vivara apresenta o melhor perfil: maior retorno esperado (23\%) e maior probabilidade de superar o CDI (46,3\%)
\end{itemize}

% ============================================================================
% 9. SCORING E OTIMIZAÇÃO
% ============================================================================
\section{Scoring Combinado e Otimização}

\subsection{Scores Z-Score por Dimensão}

\begin{center}
\begin{tabular}{lccccc}
\toprule
\textbf{Empresa} & \textbf{Return} & \textbf{Valuation} & \textbf{Quality} & \textbf{Risk} & \textbf{FINAL} \\
\midrule
Track\&Field & 0,71 & $-0,93$ & 0,54 & 0,94 & \textbf{0,31} \\
Alpargatas & $-1,17$ & $-1,18$ & $-0,26$ & $-1,33$ & $-0,98$ \\
Grupo SBF & 0,22 & 0,95 & $-1,58$ & $-0,97$ & $-0,34$ \\
Lojas Renner & $-1,26$ & 0,44 & $-0,49$ & $-0,28$ & $-0,40$ \\
\rowcolor{cinzatabela}
\textbf{Grendene} & \textbf{0,32} & \textbf{1,19} & \textbf{0,51} & \textbf{0,89} & \textcolor{verdepositivo}{\textbf{0,73}} \\
Vivara & 1,18 & $-0,48$ & 1,28 & 0,74 & \textbf{0,68} \\
\bottomrule
\end{tabular}
\end{center}

Pesos: Return 25\% + Valuation 25\% + Quality 25\% + Risk Penalty 25\%.

\subsection{Interpretação para Track\&Field}

A TFCO4 ficou em \textbf{3º lugar} no ranking com score final de 0,31:
\begin{itemize}[leftmargin=*]
    \item \textcolor{verdepositivo}{\textbf{Pontos fortes:}} Return (0,71) e Risk (0,94) --- bom retorno esperado com risco controlado
    \item \textcolor{vermelhoneg}{\textbf{Ponto fraco:}} Valuation ($-0,93$) --- múltiplos \textit{premium} penalizaram o score
    \item O mercado já precificou boa parte do crescimento; upside adicional depende de execução
\end{itemize}

\subsection{Carteira Ótima (scipy)}

A otimização contínua (máximo 40\% por ativo) resultou em:

\begin{center}
\begin{tabular}{lc}
\toprule
\textbf{Ativo} & \textbf{Peso} \\
\midrule
Grendene (GRND3.SA) & 40\% \\
Vivara (VIVA3.SA) & 40\% \\
Track\&Field (TFCO4.SA) & 20\% \\
\bottomrule
\end{tabular}
\end{center}

A Track\&Field entra na carteira ótima com 20\%, validando sua qualidade apesar do valuation esticado.

% ============================================================================
% 10. CENÁRIOS
% ============================================================================
\section{Cenários de Preço}

\begin{center}
\begin{tabular}{lcccc}
\toprule
\textbf{Cenário} & \textbf{Prob.} & \textbf{Preço-Alvo} & \textbf{Retorno} & \textbf{Premissas} \\
\midrule
\rowcolor{cinzatabela}
\textbf{Base} & 50\% & R\$ 18,57 & +12\% & SSS 8--10\%, 30--40 lojas/ano \\
\textcolor{verdepositivo}{\textbf{Bull}} & 25\% & R\$ 23,38 & +41\% & Selic 11--12\%, consumo aquecido \\
\rowcolor{cinzatabela}
\textcolor{vermelhoneg}{\textbf{Bear}} & 25\% & R\$ 11,55 & $-30\%$ & Selic 15\%+, recessão \\
\midrule
\textbf{Valor Esperado} & --- & \textbf{R\$ 18,02} & \textbf{+8,7\%} & \\
\bottomrule
\end{tabular}
\end{center}

\textbf{Cenário Base:} Crescimento regular de same-store-sales (SSS) de 8--10\%, abertura de 30--40 lojas/ano, Selic lateral. Expansão modesta de múltiplos.

\textbf{Cenário Bull:} Queda da Selic para 11--12\%, consumo aquecido, expansão de e-commerce e novos mercados. EV/EBITDA expande para 16,3x.

\textbf{Cenário Bear:} Selic sobe para 15\%+, desaceleração do consumo, pressão em margens. EV/EBITDA contrai para 10,3x.

O \textbf{valor esperado de R\$ 18,02 (+8,7\%)} é modesto e \textbf{não supera o CDI} (14,25\%), sugerindo que no cenário base o custo de oportunidade é alto.

% ============================================================================
% 11. TESTE DE HIPÓTESES
% ============================================================================
\section{Teste de Hipóteses}

\begin{center}
\begin{tabular}{lcc}
\toprule
\textbf{Hipótese} & \textbf{Resultado} & \textbf{Evidência} \\
\midrule
H1: TFCO4 tem melhor qualidade & \textcolor{verdepositivo}{\textbf{CONFIRMADA}} & ROE 28,3\% vs peers 16,1\% \\
H2: TFCO4 tem valuation atrativo & \textcolor{vermelhoneg}{\textbf{NÃO CONFIRMADA}} & EV/EBITDA 13,3x vs setor 11,6x \\
H3: TFCO4 tem menor risco & \textcolor{verdepositivo}{\textbf{CONFIRMADA}} & Vol 36,9\% vs peers 41,9\% \\
\bottomrule
\end{tabular}
\end{center}

% ============================================================================
% 12. CONCLUSÃO
% ============================================================================
\section{Conclusão e Recomendação Condicional}

\subsection{Posicionamento na Faixa de Preço}

O preço atual de R\$ 16,58 está na \textbf{posição de 80\% da faixa de 52 semanas} (R\$ 9,61 -- R\$ 18,31), indicando que está \textcolor{vermelhoneg}{\textbf{próximo do topo}}.

\subsection{Veredicto}

\begin{center}
\fcolorbox{amareloalerta}{cinzaclaro}{
\parbox{0.85\textwidth}{
\centering
\textbf{\large CUIDADO --- Preço na parte alta do range}\\[0.3cm]
Preço já andou bastante (+72\% em 12 meses). Valor esperado (+8,7\%) não supera o CDI (14,25\%). Esperar correção de 10--15\% pode ser prudente antes de iniciar posição.
}
}
\end{center}

\subsection{Estratégia Sugerida}

\begin{enumerate}[leftmargin=*]
    \item \textbf{Se preço cair para R\$ 13--14 ($<$40\% do range):} Posição cheia (3--5\% do portfólio)
    \item \textbf{Se preço ficar na faixa R\$ 14--16 (40--60\%):} Entrada parcial, acumular em correções
    \item \textbf{No nível atual R\$ 16,58 ($>$80\%):} Esperar correção de 10--15\%
    \item \textbf{Sempre:} Comparar retorno esperado vs CDI (14,25\% a.a.)
    \item \textbf{Posição máxima:} 5--8\% do portfólio (small/mid-cap)
\end{enumerate}

\subsection{Alternativas no Setor}

\begin{itemize}[leftmargin=*]
    \item \textbf{Grendene (GRND3):} Melhor score geral --- perfil defensivo com caixa líquido e dividendos altos
    \item \textbf{Vivara (VIVA3):} Segundo melhor score --- luxo acessível com margens excepcionais
    \item \textbf{Track\&Field (TFCO4):} Terceiro --- excelente qualidade, mas valuation já esticado
\end{itemize}

% ============================================================================
% LIMITAÇÕES
% ============================================================================
\section{Limitações}

\begin{itemize}[leftmargin=*]
    \item Dados fundamentalistas via yfinance podem estar desatualizados ou incompletos
    \item FCF e EBITDA dependem de disponibilidade na API
    \item Não incorpora análise de management, ESG, ou fatores qualitativos
    \item TFCO4 tem histórico curto na B3 (IPO 2020), limitando análises de longo prazo
    \item Monte Carlo assume distribuição t-Student (melhor que normal, mas ainda limitado)
    \item Betas históricos podem não refletir regime atual de Selic
    \item Sem ajuste para dividendos extraordinários ou eventos corporativos
    \item Peers podem ter modelos de negócio parcialmente diferentes (ex: Grendene = indústria)
    \item Liquidez de TFCO4 pode resultar em spreads maiores vs large-caps
    \item Dividend Yield de alguns ativos apresentou distorção na API (valores implausíveis)
\end{itemize}

\vfill
\begin{center}
\rule{0.5\textwidth}{0.4pt}\\[0.3cm]
{\small\textit{Relatório gerado automaticamente a partir de dados coletados via yfinance.\\
Análise quantitativa não substitui due diligence completa.\\
Março de 2026.}}
\end{center}

\end{document}
